\section{相关工作}

\subsection{语言模型、提示与语言智能体}

语言模型是一类经过训练、用于评估词或词元序列概率的机器学习模型。大语言模型（LLM）~\citep{gpt,gpt2,gpt3,gpt3.5,gpt4,llama,llama2}通常是指采用自回归概率分解方案、由 Transformer 架构~\citep{vaswani2017attention}参数化、拥有大量参数并在大规模语料上训练的语言模型。随着模型规模、训练数据和计算量的扩展，LLM 在生成类人文本和理解上下文方面展现出卓越能力。

另一方面，提示是释放 LLM 能力的关键。提示是控制 LLM 行为与输出的关键组件，也是人与 LLM 之间的接口。提示设计会显著影响语言模型性能；提示工程已经取得许多进展，包括上下文学习~\citep{gpt3}、思维链提示~\citep{nye2022show, wei2022chain}、ReAct~\citep{yao2022react}、自我精炼~\citep{madaan2023selfrefine}、自洽性~\citep{wang2023selfconsistency}、循环提示~\citep{zhou2023recurrentgpt}等。

语言智能体进一步扩展了语言模型的功能：它允许 LLM 使用工具~\citep{schick2023toolformer}，并将 LLM 集成到能够执行多步任务的更广泛系统中~\citep{park2023generative,hong2023metagpt,zhou2023agents,chen2023agentverse,OpenAgents}。通过把提示与工具堆叠为精心设计的流水线，智能体可以广泛应用于从客服自动化到高级数据分析的不同场景。

\subsection{从自动提示工程到智能体优化}

随着提示工程在学术界和工业界日益普及，近期许多工作研究了如何自动化提示工程过程。例如，\citet{pryzant2020automatically}和\citet{yang2024large}使用精心设计的提示，释放 LLM 自行开展提示工程的能力。另一方面，\citet{prasad-etal-2023-grips}和\citet{guo2024connecting}则采用遗传算法等不同搜索算法来优化提示。

由于提示是智能体的关键组件，自动提示工程的成功为自动智能体优化打开了可能。与自动提示工程类似，智能体优化方法也可以分为两类：\textit{基于提示}的方法和\textit{基于搜索}的方法。例如，Agent-pro~\citep{zhang2024agentpro}和 AgentOptimizer~\citep{zhang2024offline}借助精心设计的提示，优化智能体流水线某一节点中的提示或工具。这些方法作用于智能体内的孤立组件。另一条研究路线探索了基于搜索的智能体优化算法。\citet{sordoni2023joint}使用变分推断优化堆叠的 LLM；DSPy~\citep{khattab2023dspy}用搜索算法在组合空间中寻找最佳提示或节点；GPTSwarm~\citep{zhuge2024language}进一步改进了这一组合优化问题的搜索算法。这些方法存在若干主要局限。第一，只有当指标能够通过可编码的方程进行数值定义时，搜索算法才主要适用；然而，大多数智能体任务都是复杂的现实问题，其成功与否无法用某些方程定义，例如软件开发或创意写作。第二，这些方法分别更新各组件，因此会受每个节点或组件局部最优的影响。它们也不具备在流水线中添加节点或实现新工具的能力。相比之下，本文提出的\baby{}框架是首批能够“整体性”优化智能体系统的智能体学习方法之一，可以优化提示、工具、节点，以及它们堆叠成智能体的方式。

此外，近期也有一些工作通过合成数据来微调智能体的 LLM 骨干~\citep{chen2023fireact,DBLP:journals/corr/abs-2401-05268,song2024trial}。这条研究路线与本文工作正交，我们认为二者可以互补。ICE~\citep{qian2024investigateconsolidateexploit}也研究了语言智能体的跨任务迁移学习，可以与本文方法互补，用于构建自演化智能体。
